\documentclass[a4paper,11pt]{article}
\usepackage{ctex}
\usepackage{amsmath, amssymb}
\usepackage{geometry}
\geometry{left=2.5cm, right=2.5cm, top=2.5cm, bottom=2.5cm}

\begin{document}

\begin{center}
    \Large\textbf{上机作业：线性方程组的直接解法对比}
\end{center}

\section{实验目的与问题描述}
本实验旨在编写通用的线性方程组直接求解程序，包括：不选主元 Gauss 消去法、列主元 Gauss 消去法、平方根法（Cholesky 分解）以及改进的平方根法（$\text{LDL}^T$ 分解）。

利用上述程序，分别求解以下两个对称正定方程组 $Ax=b$，并对比数值解与精确解的误差：
\begin{enumerate}
    \item \textbf{100 阶三对角矩阵}：主对角线元素为 10，次对角线元素为 1，$b$ 为随机向量。
    \item \textbf{40 阶 Hilbert 矩阵}：矩阵元素 $a_{ij} = \frac{1}{i+j-1}$，右端项 $b_i = \sum_{j=1}^n \frac{1}{i+j-1}$，该构造下精确解为 $x =[1, 1, \dots, 1]^T$。
\end{enumerate}

\section{实验计算结果}
实验在双精度浮点数（Float64）环境下运行，采用无穷范数 $\|x_{\text{calc}} - x_{\text{exact}}\|_\infty$ 衡量计算误差。各算法的实际运行结果如下：

\subsection{100 阶对称正定三对角矩阵}
\begin{itemize}
    \item 不选主元 Gauss 消去法：最大误差 $= 2.0817 \times 10^{-17}$
    \item 列主元 Gauss 消去法：最大误差 $= 2.0817 \times 10^{-17}$
    \item 平方根法 (Cholesky)：最大误差 $= 2.7756 \times 10^{-17}$
    \item 改进的平方根法 ($\text{LDL}^T$)：最大误差 $= 1.3878 \times 10^{-17}$
\end{itemize}

\subsection{40 阶 Hilbert 矩阵}
\begin{itemize}
    \item 不选主元 Gauss 消去法：最大误差 $= 1.4891 \times 10^2$
    \item 列主元 Gauss 消去法：最大误差 $= 5.1932 \times 10^2$
    \item 平方根法 (Cholesky)：程序抛出异常中止（矩阵分解过程中出现对角元小于等于0，触发负数开方错误）
    \item 改进的平方根法 ($\text{LDL}^T$)：最大误差 $= 5.9162 \times 10^1$
\end{itemize}

\section{结果分析与讨论}

基于上述实验数据，对直接解法在不同条件下的表现分析如下：

\subsection{良态方程组的求解与对称性的利用}
100 阶三对角矩阵属于严格对角占优的对称正定矩阵，条件数较小，是一个典型的良态问题。
由实验结果可知，四种方法均达到了 $10^{-17}$ 量级的计算精度，接近机器精度极限。这表明对于良态矩阵，即使不使用列主元策略，也不会引发明显的舍入误差放大。

从算法效率来看，平方根法和改进的平方根法利用了矩阵的对称性，只需处理下三角部分，其计算量和存储量约为 Gauss 消去法的一半。其中，$\text{LDL}^T$ 分解因为避免了开方运算，运算效率更高，且在本组测试中表现出了最优的计算精度。

\subsection{病态方程组对直接解法的影响}
40 阶 Hilbert 矩阵是典型的病态矩阵，其条件数随阶数增加呈指数级增长。
在第二组实验中，所有算法的误差均达到了 $10^1 \sim 10^2$ 量级（而真实解均为 1），说明有效数字已经全部丢失，计算结果完全失真。

实验中观察到一个典型现象：平方根法（Cholesky 分解）在执行过程中发生了程序中断。这是因为在有限精度的浮点数运算中，舍入误差被 Hilbert 矩阵的病态性剧烈放大，导致中间某一步骤的对角元出现了非正数。这破坏了理论上的数值正定性，使得后续的开方运算无法进行。相比之下，$\text{LDL}^T$ 分解由于全程没有开方运算，虽然结果同样存在巨大的误差，但程序能够执行完毕，表现出了相对较好的鲁棒性。

\subsection{结论}
数值稳定的算法（如列主元 Gauss 消去法）能够有效控制算法执行步骤中产生的误差累积，但无法消除由矩阵固有病态性带来的误差放大。对于对称正定且良态的稠密方程组，$\text{LDL}^T$ 分解通常是兼顾效率和稳定性的首选直接解法；而面对高阶 Hilbert 矩阵这类病态方程组，常规的直接解法往往失效，在实际工程中需考虑采用多精度数据类型或引入正则化等其他数值手段。

\end{document}